Existing game evaluation frameworks measure quality rather than design character, producing
total scores that obscure how games differ from one another.
This paper formalizes the five empirically-discovered dimensions from TG-A2 into the
TIGER framework: Tension (T), Interaction (I), Game-Architecture (G), Elegance (E), and Range (R).
Each dimension is defined with anchor descriptors at 0, 5, and 10 derived from the
150-game TIGRIS corpus (TG-A1) and instantiated as concrete TIGER profiles for 20
canonical games spanning BGG weight 1.07--4.09.
We validate TIGER via a forced-choice triplet study with five expert raters
(game designer, theorist, hobbyist, publisher, journalist) and 19 triplets,
achieving \textbf{98.9\% pairwise discrimination accuracy} (94/95 judgments;
mean rater confidence 4.57/5), far exceeding the pre-registered 90\% threshold.
18 of 19 triplets were unanimous (5-0); the single split case is analytically
coherent and reveals a legitimate single-dimension comparison scenario.
The central argument is that TIGER measures \emph{distinctiveness} rather than quality:
two games can share an identical BGG rating while occupying opposite corners of
TIGER space --- Crokinole (T=3.6, G=2.4) vs.\ Agricola (T=8.2, G=7.6) differ by
8.2 TIGER units despite similar player engagement.
TIGER is useful precisely where quality frameworks fail: for recommendation,
design diagnosis, and identifying unexplored design regions.
